\section{An Egyptian-fraction recursion for \texorpdfstring{\seq{A003167}}{A003167}}
\label{sec:a003167}

An integral \(n\)-dimensional cuboid with nondecreasing side lengths
\(x_1\le\cdots\le x_n\) has volume
\(\prod_i x_i\).  Its total facet hyperarea is
\(2\sum_i\prod_{j\ne i}x_j\).  Equality of volume and total facet hyperarea is
therefore equivalent to
\begin{equation}
 \frac1{x_1}+\cdots+\frac1{x_n}=\frac12.
 \label{eq:a003167-egypt}
\end{equation}
Thus \seq{A003167} counts nondecreasing positive-integer solutions of
\cref{eq:a003167-egypt}.

For coprime integers \(p\ge0\), \(q\ge1\), an integer \(r\ge0\), and a lower
bound \(L\ge1\), let
\[
 C(r,p,q,L)
\]
be the number of nondecreasing \(r\)-tuples
\(L\le x_1\le\cdots\le x_r\) whose reciprocal sum is \(p/q\).
Use the terminal conventions
\[
 C(0,0,1,L)=1,
 \qquad
 C(0,p,q,L)=0\quad(p>0),
\]
and \(C(r,0,1,L)=0\) for \(r>0\).  For \(r=1\), the value is one exactly
when \(p>0\), \(p\mid q\), and \(q/p\ge L\).

\begin{theorem}[Exact finite recursion]
\label{thm:a003167-recursion}
For \(r\ge2\) and \(p>0\),
\begin{equation}
 C(r,p,q,L)=
 \sum_{x=\max\{L,\lfloor q/p\rfloor+1\}}^{\lfloor rq/p\rfloor}
 C(r-1,p_x,q_x,x),
 \label{eq:a003167-recursion}
\end{equation}
where \(p_x/q_x\) is the reduced form of
\begin{equation}
 \frac{px-q}{qx}.
 \label{eq:a003167-residual}
\end{equation}
In particular,
\begin{equation}
 \boxed{\seq{A003167}(n)=C(n,1,2,1)}.
 \label{eq:a003167-main}
\end{equation}
\end{theorem}

\begin{proof}
Suppose \(x=x_1\) is the first denominator.  Since \(r-1\) positive
reciprocals remain, the residual sum must be positive, so
\(1/x<p/q\) and therefore \(x>q/p\).  Monotonicity also requires \(x\ge L\).
On the other hand every reciprocal is at most \(1/x\), whence
\(p/q\le r/x\), or \(x\le rq/p\).  These are exactly the finite bounds in
\cref{eq:a003167-recursion}.  After choosing \(x\), subtracting \(1/x\)
produces \cref{eq:a003167-residual}, and the remaining denominators are at
least \(x\).  Every solution has a unique first denominator and every summand
constructs valid solutions, proving the recursion.  Equation
\cref{eq:a003167-main} is the reformulation
\cref{eq:a003167-egypt}.
\end{proof}

\subsection{Two-denominator divisor terminal}

When two denominators remain, the equation
\[
 \frac1x+\frac1y=\frac pq
\]
is equivalent to
\begin{equation}
 (px-q)(py-q)=q^2.
 \label{eq:a003167-factor}
\end{equation}
Hence the terminal count can be obtained by enumerating divisors \(d\mid q^2\)
and testing
\begin{equation}
 x=\frac{d+q}{p},
 \qquad
 y=\frac{q^2/d+q}{p},
 \qquad
 L\le x\le y.
 \label{eq:a003167-terminal}
\end{equation}
This is an exact acceleration of \cref{eq:a003167-recursion}, not an added
assumption.  The implementation uses deterministic integer factorization and
protects the products in \(q^2\) with widened integer arithmetic.

\subsection{The seventh term}

\begin{resultbox}[Exact evaluation]
\[
 \boxed{\seq{A003167}(7)=155068098.}
\]
\end{resultbox}

For \(n=7\), the first denominator lies in \(3\le x_1\le14\).  The expensive
branch \(x_1=3\) was refined into the following disjoint pieces, each dictated
by the same denominator bounds:
\begin{align*}
 &x_1=4,\ldots,14;\\
 &x_1=3,\quad x_2=8,\ldots,36;\\
 &x_1=3,\quad x_2=7,\quad x_3=44,\ldots,210;\\
 &x_1=3,\quad x_2=7,\quad x_3=43,
      \quad x_4=1807,\ldots,7224,
\end{align*}
with the last interval partitioned into consecutive subranges.  The
aggregation checker verifies the embedded prefixes and ranges, proves that the
\(x_4\)-subranges cover \([1807,7224]\) without gaps or overlaps, and sums 262
branch counts.  Altogether the branches visit 151379115 recursion nodes and
make 151271985 terminal pair calls.

The same C++ recursion reproduces the known values
\[
 2,\ 10,\ 108,\ 2892,\ 270332
\]
for \(n=2,\ldots,6\).  A separate Python implementation using exact rational
recursion, without the factorization terminal, independently reproduces this
prefix.  The theorem is all-dimensional; the branch computation supplies the
new seventh numerical value.
